\documentclass[sigconf,natbib]{acmart}

\usepackage{booktabs}
\usepackage{xcolor}
\usepackage{amsmath}
\usepackage{amssymb}
\usepackage{listings}

%% ACM/ISWC conference metadata
\acmConference[ISWC '25]{International Semantic Web Conference}{November 2025}{Online}
\acmYear{2025}
\copyrightyear{2025}
\setcopyright{none}

\begin{document}

\title{Auditing Deliberative Alignment in EU Lawmaking:\\
The Deliberation Knowledge Graph and a Pilot Metric Study}

\author{Simone Vagnoni}
\affiliation{%
  \institution{CIRSFID -- Alma Mater Studiorum, Universit\`a di Bologna}
  \city{Bologna}
  \country{Italy}
}
\email{simone.vagnoni3@unibo.it}

\author{V\'ictor Rodr\'iguez-Doncel}
\affiliation{%
  \institution{Ontology Engineering Group, Universidad Polit\'ecnica de Madrid}
  \city{Madrid}
  \country{Spain}
}
\email{vrodriguez@fi.upm.es}

\begin{abstract}
Assessing whether citizen input to EU public consultations actually shapes legislative deliberation remains an open empirical challenge.
Existing studies cluster consultation responses or analyse parliamentary speeches in isolation; few attempt systematic cross-forum comparison, and none do so through a reusable semantic infrastructure.
This paper applies the Deliberation Knowledge Graph (DKG)~\cite{Vagnoni2026}---a previously published OWL ontology and Linked Open Data infrastructure linking EU public consultations (\textit{Have Your Say}) to European Parliament plenary debates via EuroLex procedure identifiers---to a novel analytical problem: quantifying deliberative distance between citizens and legislators across multiple legislative dossiers.

We introduce the \textit{Deliberative Distance Index} (DDI), a composite metric computed directly over the DKG's bipartite contribution--topic subgraph.
The DDI combines three dimensions: topic distance (avg-max cosine similarity between multilingual BERTopic centroids, range $[0,1]$), salience distance (Jensen--Shannon divergence log$_2$ of topic-weight distributions, range $[0,1]$), and stance distance (JSD log$_2$ of zero-shot NLI-classified stances, range $[0,1]$).
We further derive topological graph metrics---cross-phase bridging ratio (CPBR), mean phase entropy (MPET), and cross-phase degree correlation ($r_{BC}$)---from the graph structure alone, without relying on language models.
The full pipeline operates natively across 24 EU languages without translation, using a multilingual sentence encoder and the mDeBERTa-v3 NLI model.

Applied to a pilot study of three EU legislative dossiers, the framework reveals that thematic salience divergence is the dominant driver of cross-forum distance.
Bootstrap and permutation tests yield a differentiated picture: post-consultation $\Delta$DDI is not statistically different from zero for COVID and Climate Target ($p > 0.6$), while for the Pesticides SUR dossier $\Delta\text{DDI} = -0.027$ is statistically significant ($p = 0.028$, bootstrap CI $[-0.046, -0.006]$), indicating post-consultation divergence consistent with the SUR regulation's withdrawal in February 2024.
These results demonstrate that the DKG enables reproducible, structured auditing of deliberative alignment at a scale and multilingual coverage previously infeasible.
Code and SPARQL endpoint are publicly available.
\end{abstract}

\keywords{deliberative democracy, knowledge graph, public consultations, EU lawmaking, topic modeling, stance detection, deliberative distance, Linked Open Data, FAIR data, computational social science}

\maketitle

\section{Introduction}

\subsection{Motivation}
Public consultations are central to EU policymaking under the Better Regulation agenda \citep{EC2015}, yet empirical evidence on whether citizen feedback influences legislative deliberation remains scarce. While consultations generate hundreds of thousands of responses, systematic analysis of their alignment with parliamentary debates has been limited by:
\begin{itemize}
  \item \textbf{Data fragmentation}: Consultation responses, parliamentary transcripts, and legal texts exist in siloed formats.
  \item \textbf{Methodological gaps}: Existing studies cluster consultation responses \citep{DiPorto2024} but rarely link them to legislative outcomes.
  \item \textbf{Scale challenges}: Manual traceability analysis does not scale to thousands of consultations and debates.
\end{itemize}

This paper applies the DKG~\cite{Vagnoni2026} to a novel analytical problem, contributing:
(1) a composite metric---the Deliberative Distance Index---computed over the graph's bipartite structure;
(2) a multilingual NLP pipeline operating across 24 EU languages without translation;
(3) topological graph metrics derived from the contribution--topic subgraph; and
(4) empirical findings across three EU legislative dossiers.

\subsection{Research Questions}
\begin{enumerate}
  \item[\textbf{RQ1}] Do EU public consultations and parliamentary debates discuss the same topics, and with similar thematic emphasis?
  \item[\textbf{RQ2}] When topics overlap, do citizens and legislators hold similar stances toward the policy proposition?
  \item[\textbf{RQ3}] Does post-consultation parliamentary debate show measurable convergence toward citizen concerns?
  \item[\textbf{RQ4}] Which structural properties of the contribution--topic graph predict deliberative alignment?
\end{enumerate}

\subsection{Contributions}
\begin{itemize}
  \item[\textbf{C1}] \textbf{Deliberative Distance Index (DDI)}: A composite metric with three interpretable dimensions---topic distance, salience distance, and stance distance---computed via symmetric avg-max cosine similarity and Jensen--Shannon divergence over the DKG's bipartite subgraph.
  \item[\textbf{C2}] \textbf{Topological graph metrics}: Cross-phase bridging ratio (CPBR), mean phase entropy (MPET), and cross-phase degree correlation ($r_{BC}$), derived from graph structure alone without language models.
  \item[\textbf{C3}] \textbf{Multilingual pipeline}: BERTopic with multilingual sentence embeddings plus mDeBERTa-v3 zero-shot NLI, operating natively across 24 EU languages.
  \item[\textbf{C4}] \textbf{Pilot study}: Three EU legislative dossiers (EU Digital COVID Certificate, EU 2040 Climate Target, Sustainable Use of Pesticides Regulation), covering citizen consultation feedback matched to European Parliament plenary debates via temporal stratification.
  \item[\textbf{C5}] \textbf{Open pipeline}: All scripts, intermediate data, and the SPARQL endpoint are publicly available for replication and extension.
\end{itemize}

\section{Related Work}

\subsection{Computational Analysis of Public Consultations}
\citet{DiPorto2024} apply NLP clustering to 830 EU consultation responses, identifying topic distributions but not linking to legislative outcomes.
\citet{Bunea2024} analyse 42 EU Commission consultations and find that over 70\% present low-dimensional policy spaces (one or two dominant dimensions), suggesting that thematic diversity in citizen input is structurally constrained.
\citet{Neumann2025} provide empirical evidence that HYS is used predominantly by organised stakeholders rather than ordinary citizens, owing to low participatory literacy among the EU public---contextualising why consultation-to-parliament alignment may be limited in practice.
Our work extends these by matching consultation corpora to parliamentary debates and measuring alignment via the DDI.

\subsection{Parliamentary Debate Analysis}
Studies analyze EP speeches \citep{Abercrombie2020} and voting behavior \citep{Hix2006}, but rarely connect to pre-legislative consultation. We bridge this gap by pairing consultation corpora with debate transcripts.

\subsection{Deliberative Quality Metrics}
Deliberative quality has been measured through manual coding \citep{Steenbergen2003} or argument mining \citep{Lawrence2020}. We propose a computational metric (DDI) for cross-forum comparison at scale.

\subsection{Cross-Corpus Topic Comparison}
\citet{Adam2024} independently introduce Bidirectional Topic Matching (BTM), a framework that quantifies thematic overlap between corpora by training separate topic models and applying them reciprocally via co-occurrence-based pairing strength.
Our $d_{\text{topic}}$ uses a different but related approach---symmetric avg-max cosine similarity between per-phase BERTopic centroids---which Adam \& Kogler validate as a strong baseline for BTM (Cohen's $\kappa = 0.75$--$0.81$).
The two approaches are independently derived and complementary: BTM emphasises co-occurrence structure; our centroid-based method emphasises geometric distance in embedding space.

\subsection{Knowledge Graphs for Civic Data}
Semantic web technologies have been applied to legal texts \citep{Palmirani2018} and legislative procedures \citep{Mader2012}, but rarely to multi-platform deliberation. The DKG extends this to heterogeneous civic participation data.

\section{Deliberation Knowledge Graph Architecture}

\subsection{Motivation for a Unified Ontology}
Deliberative data exists across platforms (consultations, forums, debates, court arguments) with incompatible schemas. A knowledge graph provides:
\begin{itemize}
  \item \textbf{Semantic interoperability}: Shared vocabulary for contributions, topics, actors, stances.
  \item \textbf{Traceability}: Links from consultation feedback → EP debate → legal text.
  \item \textbf{SPARQL queries}: Structured retrieval across platforms (e.g., "Find all AI-related contributions by NGOs").
\end{itemize}

\subsection{Ontology Design}
The Deliberation Ontology (\url{https://w3id.org/deliberation/ontology}) defines:

\textbf{Core Classes}:
\begin{itemize}
  \item \texttt{del:DeliberationProcess} (e.g., HYS consultation, EP debate session)
  \item \texttt{del:Contribution} (feedback, speech, forum post)
  \item \texttt{del:Participant} (citizen, MEP, stakeholder)
  \item \texttt{del:Topic} (extracted via BERTopic or manual tagging)
  \item \texttt{del:Stance} (support, oppose, neutral)
  \item \texttt{del:Argument} (claim + evidence structure)
\end{itemize}

\textbf{Properties}:
\begin{itemize}
  \item \texttt{del:madeBy} (links contribution to participant)
  \item \texttt{del:hasTopic} (links to topic entity)
  \item \texttt{del:hasStance} (stance classification)
  \item \texttt{del:isResponseTo} (links consultation to legislative procedure via CELEX/Ares references)
\end{itemize}

\subsection{Data Sources Integrated}

The DKG~\cite{Vagnoni2026} integrates heterogeneous deliberative sources under the unified \texttt{del:} ontology.
This paper uses the HYS and EP Plenary subgraphs; other sources modelled in the DKG (online deliberation platforms, structured debate corpora) are left for future pilots.

\begin{table}[h]
\centering
\small
\begin{tabular}{llll}
\toprule
\textbf{Source} & \textbf{Period} & \textbf{Type} & \textbf{Used in pilot} \\
\midrule
EU Have Your Say (HYS)         & 2016--2025 & Citizen consultation feedback        & Yes \\
EP Plenary Debates             & 2021--2025 & Parliamentary speeches (verbatim)    & Yes \\
\midrule
\multicolumn{4}{l}{\textit{Modelled in DKG ontology; not used in this pilot:}} \\
Decide Madrid (Consul)         & 2015--2019 & Participatory democracy proposals    & No \\
Decidim Barcelona              & 2016--2022 & Online participatory platform posts  & No \\
DeliData                       & 2019--2021 & Multi-party deliberation transcripts & No \\
Pol.is (vTaiwan, EU AI Panel)  & 2021--2023 & Large-scale opinion matrix data      & No \\
SCOTUS Oral Arguments          & 2017--2021 & Legal deliberation transcripts       & No \\
\bottomrule
\end{tabular}
\caption{Data sources modelled in the DKG ontology (\url{https://w3id.org/deliberation/ontology}).
This paper uses the HYS and EP Plenary subgraphs; other sources are available in the public endpoint.}
\label{tab:sources}
\end{table}

\subsection{RDF Triple Store and SPARQL Endpoint}
The DKG is served via Apache Jena Fuseki with:
\begin{itemize}
  \item \textbf{Named graphs} per source (e.g., \texttt{:HaveYourSay}, \texttt{:EPDebates})
  \item \textbf{Federated queries} across platforms
  \item \textbf{OWL reasoning} for inferring topic hierarchies
\end{itemize}

Example SPARQL query:
\begin{lstlisting}[language=SQL, basicstyle=\small\ttfamily]
PREFIX del: <https://w3id.org/deliberation/ontology#>

SELECT ?contribution ?text ?topic WHERE {
  ?contribution a del:Contribution ;
                del:hasTopic ?topic ;
                del:text ?text .
  ?topic del:name "artificial intelligence" .
  FILTER (CONTAINS(LCASE(?text), "ethical"))
}
\end{lstlisting}

\subsection{Comparison to Existing Legal KGs}
Unlike EUR-Lex (legal texts only) or OEIL (procedures only), the DKG integrates \textit{deliberative processes} leading to legislation, enabling traceability from citizen input to legal outcome.

\section{Pilot Study: HYS-EP Comparison}

\subsection{Case Selection Rationale}

We select three EU legislative dossiers for which all three data strata are available:
a Phase~B citizen corpus from HYS, Phase~A EP speeches prior to consultation close, and
Phase~C EP speeches after. An additional criterion was applied: no single language may exceed
45\% of Phase~B responses, to exclude cases dominated by coordinated single-country campaigns,
a known distortion in HYS data.

The three cases provide deliberate variation across two axes:

\begin{itemize}
  \item \textbf{Corpus size and participation profile}: COVID Certificate attracted mass
        anonymous individual response (5,000 sampled from 296,600 submissions); Climate
        Target attracted moderate, ideologically diverse response (568 items);
        Pesticides attracted expert-heavy stakeholder response (8,140 items).
  \item \textbf{Policy domain and stance polarity}: COVID is high-salience and politically
        polarised; Climate Target is ideologically contested across the pro/against
        climate-ambition axis; Pesticides is expertise-driven, opposing agricultural
        industry and environmental NGOs.
\end{itemize}

\subsection{Data Collection}

HYS feedback is extracted from a local SQLite mirror of the EC Better Regulation portal
(scraped 2024--2025 via the public API). EP speeches are drawn from a verbatim JSON corpus
covering 2021--2025 EP plenary sessions, filtered by dossier-specific keywords and partitioned
temporally relative to the HYS consultation window.

\begin{table*}[t]
\centering
\small
\begin{tabular}{llrrrl}
\toprule
\textbf{Case} & \textbf{HYS ref} & \textbf{Ph.~A} & \textbf{Ph.~B} & \textbf{Ph.~C} & \textbf{Window} \\
\midrule
EU Digital COVID Certificate  & COM(2022)50       & 419   & 4,624  & 86    & Mar--May 2022 \\
EU 2040 Climate Target        & Ares(2023)2344532 & 2,951 & 568    & 1,234 & Apr--Jun 2023 \\
Sustainable Use of Pesticides & Ares(2020)2804518 & 388   & 8,140  & 800   & Jun--Sep 2022 \\
\bottomrule
\end{tabular}
\caption{Pilot study corpus. Ph.~A/C = EP speeches before/after consultation close.
Ph.~B = HYS items with free text ($>$30 chars). COVID Phase~B is a random sample (seed~42)
from 296,600 total submissions.}
\label{tab:dataset}
\end{table*}

\noindent
All data derive from public EU sources and are released under CC~BY~4.0.
RDF exports (Turtle) are available in the DKG SPARQL endpoint.

\section{Methodology}

\subsection{Overview of the Phase-Stratified Pipeline}

A key methodological challenge in measuring deliberative distance is temporal: EU public
consultations may open before, during, or after parliamentary deliberation on the same dossier.
Treating the two corpora as synchronous would conflate pre-existing legislative discourse with
any post-consultation convergence. We address this through a \textit{phase-stratified} design
that partitions EP speeches according to their temporal relation to the consultation window,
enabling a quasi-experimental pre/post comparison.

The full pipeline proceeds in six stages:

\begin{enumerate}
  \item \textbf{Dossier anchoring}: Map each case to its EUR-Lex procedure number; retrieve
        all EP verbatim records linked to that procedure via the EP Open Data Portal.
  \item \textbf{Temporal stratification}: Partition EP speeches into Phase~A (before
        consultation close) and Phase~C (after consultation close); the HYS corpus forms
        Phase~B.
  \item \textbf{Topic modeling}: Run BERTopic separately on each phase for $d_{\text{topic}}$
        centroids; also fit a unified BERTopic model on all phases jointly (Step~1b) to provide
        the shared topic vocabulary required for $d_{\text{salience}}$.
  \item \textbf{Cross-forum alignment}: Compute cosine similarity between topic embeddings
        across corpora; classify topic gaps and aligned pairs.
  \item \textbf{Stance detection}: Apply zero-shot NLI with a dossier-level proposition
        to each document; aggregate per-forum stance distributions.
  \item \textbf{DDI computation and validation}: Compute $\text{DDI}(B, A)$ and
        $\text{DDI}(B, C)$; derive the \textit{consultation effect} $\Delta\text{DDI}$;
        perform sensitivity analysis over weight configurations.
\end{enumerate}

\subsection{Step 1: Dossier Anchoring via EUR-Lex}

Keyword-based retrieval of EP speeches introduces noise from debates that mention a topic
incidentally. We instead anchor each case to its official EUR-Lex procedure identifier
keyword-based filtering on the EP verbatim corpus (2021--2025), then partition by
consultation close date (see Table~\ref{tab:dataset}).

\begin{table}[h]
\centering
\small
\begin{tabular}{llll}
\toprule
\textbf{Case} & \textbf{HYS reference} & \textbf{EUR-Lex procedure} & \textbf{Consultation window} \\
\midrule
EU Digital COVID Certificate  & COM(2022)50       & 2022/0031(COD) & Mar--May 2022 \\
EU 2040 Climate Target        & Ares(2023)2344532 & COM(2023)653   & Apr--Jun 2023 \\
Sustainable Use of Pesticides & Ares(2020)2804518 & 2022/0196(COD) & Jun--Sep 2022 \\
\bottomrule
\end{tabular}
\caption{Case anchoring: HYS references, EUR-Lex procedure identifiers, and consultation windows.}
\label{tab:cases}
\end{table}

\subsection{Step 2: Temporal Stratification}

Let $t_{\text{open}}$ and $t_{\text{close}}$ denote the opening and closing dates of the
HYS consultation. We define three temporal strata:

\begin{align}
\text{Phase A (EP-pre):}   &\quad \{s \in \mathcal{S}_E \mid \text{date}(s) < t_{\text{close}}\} \\
\text{Phase B (HYS):}      &\quad \{f \in \mathcal{F}_H \mid t_{\text{open}} \leq \text{date}(f) \leq t_{\text{close}}\} \\
\text{Phase C (EP-post):}  &\quad \{s \in \mathcal{S}_E \mid \text{date}(s) > t_{\text{close}}\}
\end{align}

where $\mathcal{S}_E$ is the set of EP speeches for the dossier and $\mathcal{F}_H$ is the
HYS feedback corpus.

\noindent\textbf{Note on Phase~A temporal bounds.}
Using $t_{\text{close}}$ as the Phase~A cutoff means that speeches delivered \emph{during}
the consultation window ($t_{\text{open}} \leq \text{date}(s) < t_{\text{close}}$) are
included in Phase~A. To quantify this overlap, we inspect the actual date ranges of Phase~A
speeches: for COVID, 29.1\% of Phase~A speeches fall within the consultation window
(122 of 419); for Climate Target, 37.3\% (1,102 of 2,951); for Pesticides, only 0.8\%
(3 of 388). Using $t_{\text{open}}$ as the cutoff (``Phase~A\textsubscript{strict}'')
is methodologically cleaner; as a robustness check we recompute $d_{\text{salience}}$ and
$d_{\text{stance}}$ for Phase~A\textsubscript{strict} and verify that DDI(B, A) changes by
at most 0.012 (COVID) and 0.004 (Climate), without altering qualitative conclusions
(Section~\ref{sec:robustness}). The strict definition is recommended for future work.

\subsection{Step 3: Topic Modeling with BERTopic}

\subsubsection{Multilingual Embedding}

HYS feedback arrives in up to 24 EU official languages. We embed all documents using
\texttt{paraphrase-multilingual-MiniLM-L12-v2} (384 dimensions), which projects texts from
different languages into a shared semantic space without prior translation. EP speeches are
predominantly English; the same model is applied for consistency.

\subsubsection{Model Architecture}
\begin{itemize}
  \item \textbf{Dimensionality reduction}: UMAP (384 $\rightarrow$ 5 dimensions, cosine metric,
        \texttt{n\_neighbors}=15)
  \item \textbf{Clustering}: HDBSCAN (\texttt{min\_cluster\_size}=10 for Phase~A/C,
        \texttt{min\_cluster\_size}=50 for Phase~B given larger corpus)
  \item \textbf{Keyword extraction}: class-based TF-IDF (c-TF-IDF)
  \item \textbf{Large corpora}: For Phase~B with $>$50k documents (COVID), we apply
        online BERTopic with mini-batch updates to manage memory.
\end{itemize}

\subsubsection{Output per Stratum}
For each stratum and case:
\begin{itemize}
  \item Topic list $\{T^k\}$ with top-10 keywords and representative documents
  \item Topic prevalence $p^k = |D(T^k)| / |D|$
  \item Centroid embedding $\mathbf{c}^k \in \mathbb{R}^{384}$ (mean of member document embeddings)
\end{itemize}

\subsection{Step 4: Cross-Forum Topic Alignment}

\subsubsection{Similarity Matrix}
Given HYS topics $\{T_H^i\}$ and EP topics $\{T_E^j\}$ (from Phase~A or Phase~C):
\begin{equation}
S_{ij} = \frac{\mathbf{c}_H^i \cdot \mathbf{c}_E^j}{\|\mathbf{c}_H^i\| \, \|\mathbf{c}_E^j\|}
\end{equation}

\subsubsection{Matching Rule}
Topics $T_H^i$ and $T_E^j$ are \textit{aligned} if:
\begin{equation}
S_{ij} = \max_k S_{ik} \quad \text{and} \quad S_{ij} > \theta
\end{equation}
where $\theta = 0.65$ is a conservative threshold consistent with high-similarity ranges
for multilingual sentence embeddings \citep[cf.][]{Reimers2019}.
Note that $\theta$ affects topic \textit{classification} (aligned/partial/gap) but not the
DDI computation itself, which is based on the continuous avg-max cosine score over all topics.

\subsubsection{Gap Classification}
\begin{itemize}
  \item \textbf{Civic-only}: $\max_j S_{ij} < 0.5$ — topics raised by citizens absent from
        parliamentary discourse (potential representation gaps)
  \item \textbf{Institutional-only}: $\max_i S_{ij} < 0.5$ — technical/procedural topics
        absent from public consultation
  \item \textbf{Aligned}: $S_{ij} \geq \theta$ — topics present in both forums
\end{itemize}

\subsection{Step 5: Stance Detection}

\subsubsection{Proposition Design}
We assign one normatively grounded policy proposition per dossier, representing the core
regulatory question at stake:

\begin{itemize}
  \item \textbf{COVID}: ``The EU digital COVID certificate should be extended.''
  \item \textbf{Climate Target}: ``The EU should set ambitious climate targets for 2040 to achieve carbon neutrality by 2050.''
  \item \textbf{Pesticides}: ``The use of pesticides in EU agriculture should be reduced.''
\end{itemize}

Using a single, dossier-level proposition rather than per-topic propositions has two
advantages: (1) it makes stance distributions directly comparable across phases; (2) it
avoids label inconsistency that would arise from topic-specific propositions differing
in framing. The trade-off is that sub-topic nuance is collapsed into one aggregate dimension.

\subsubsection{Zero-Shot NLI Classification}
We apply \texttt{MoritzLaurer/mDeBERTa-v3-base-xnli-multilingual-nli-2mil7}, a multilingual
NLI model supporting all 24 EU official languages, avoiding the translation bottleneck of
English-only models such as BART-MNLI. Each text is classified as \textit{entailment}
(support), \textit{neutral}, or \textit{contradiction} (oppose) relative to the proposition:

\begin{lstlisting}[language=Python, basicstyle=\small\ttfamily]
from transformers import pipeline
nli = pipeline("zero-shot-classification",
    model="MoritzLaurer/mDeBERTa-v3-base-xnli-multilingual-nli-2mil7")
result = nli(text, candidate_labels=[proposition],
             hypothesis_template="{}")  # proposition used directly as hypothesis
\end{lstlisting}

\noindent
We pass the proposition directly as the NLI hypothesis (via \texttt{hypothesis\_template="\{\}"}
rather than a wrapping template such as \texttt{"This text is about \{\}"}),
since the proposition is already a full declarative statement designed for NLI entailment.
This follows the usage pattern recommended for mDeBERTa-v3 when the label is a complete sentence.

\subsubsection{Stance Distribution Comparison}
Stance labels are aggregated per phase to produce corpus-level distributions
$(p_{\text{support}}, p_{\text{neutral}}, p_{\text{oppose}})$.
Cross-phase divergence is then measured as Jensen--Shannon divergence (see $d_{\text{stance}}$
in Section~\ref{sec:ddi}).

\subsection{Step 6: Deliberative Distance Index (DDI)}
\label{sec:ddi}

\subsubsection{Component Metrics}

The DDI is computed at the corpus (phase) level using three distance components between
Phase~B (citizens) and Phase~$X \in \{A, C\}$ (EP):

\paragraph{Topic Distance} ($d_{\text{topic}}$) — semantic divergence between the two
topic sets, measured as the symmetric average-maximum cosine similarity between BERTopic
centroids:
\begin{equation}
d_{\text{topic}}(B, X) = 1 - \frac{1}{2}\left(
  \frac{1}{|T_B|}\sum_{i}\max_j S_{ij}
  + \frac{1}{|T_X|}\sum_j \max_i S_{ij}
\right)
\end{equation}
where $S_{ij} = \cos(\mathbf{c}_B^i, \mathbf{c}_X^j)$ and $\mathbf{c}^k$ is the centroid
of topic $k$. The symmetric formulation avoids the topic-count asymmetry bias of Jaccard.

\paragraph{Salience Distance} ($d_{\text{salience}}$) — divergence in thematic emphasis,
measured as the Jensen--Shannon divergence \citep{Lin1991} with log base~2 (range $[0,1]$)
between corpus-level topic-weight distributions:
\begin{equation}
d_{\text{salience}}(B, X) = \mathrm{JSD}_2\!\left(\{p_B^k\}_k \;\|\; \{p_X^k\}_k\right) \in [0,1]
\end{equation}
where $p^k = |D(T^k)| / |D|$ is the fraction of documents assigned to topic $k$.
We use the unweighted (symmetric) JSD rather than a corpus-size-weighted variant;
\citet{Lu2020} show that size-weighted JSD produces misleading results when the corpora
compared differ substantially in size, as is the case here (e.g.\ COVID Phase~B: 4,624 items total,
of which 3,097 assigned to topics vs.\ Phase~C: 86 items, of which 80 assigned).
Crucially, both distributions are defined over the \emph{unified} BERTopic vocabulary---a
single model fitted jointly on all phases (Step~1b of the pipeline)---so that index $k$
refers to the same semantic cluster in $p_B$ and $p_X$.
Computing JSD over per-phase topic IDs would be semantically invalid because each run
assigns IDs independently (topic ``0'' = largest cluster in that run), making cross-phase
ID comparison meaningless.
By contrast, $d_{\text{topic}}$ uses \emph{per-phase} BERTopic centroids: since avg-max cosine
finds the geometrically nearest centroid pair regardless of ID labels, it does not require
a shared topic index---only that centroids lie in the same embedding space (the shared
384-dim paraphrase-multilingual-MiniLM-L12-v2 space satisfies this).

\paragraph{Stance Divergence} ($d_{\text{stance}}$) — divergence in aggregate policy
positions, measured as JSD (log base~2) between corpus-level stance distributions:
\begin{equation}
d_{\text{stance}}(B, X) = \mathrm{JSD}_2(P_B^{\text{stance}} \| P_X^{\text{stance}}) \in [0, 1]
\end{equation}
where $P^{\text{stance}} = (p_{\text{support}}, p_{\text{neutral}}, p_{\text{oppose}})$
is the aggregate distribution produced by the NLI classifier applied to all documents
in the respective phase.

\subsubsection{DDI Formula}

The DDI aggregates three distance components at the corpus (phase) level:

\begin{equation}
\text{DDI}(B, X) = w_t \cdot d_{\text{topic}}(B,X)
                 + w_s \cdot d_{\text{salience}}(B,X)
                 + w_\sigma \cdot d_{\text{stance}}(B,X)
\end{equation}

\noindent All three components are computed at corpus level (not per aligned-topic pair),
so $\text{DDI}(B, X)$ is a single scalar per phase comparison.
Default weights: $w_t=0.40$, $w_s=0.35$, $w_\sigma=0.25$ (sum to 1).
These weights are set \emph{a priori} based on a theoretical judgment that topic coverage
constitutes the primary dimension of deliberative alignment; they are not empirically
calibrated. We treat them as a design choice requiring robustness testing rather than as
a claim about the relative importance of the three dimensions \citep[cf.][]{OECD2008}.
Consultation inclusivity (Shannon entropy of Phase~B actor types) is computed as a
descriptive variable only and not included in the DDI, since it measures a property of
Phase~B alone rather than cross-phase distance.

\subsubsection{Consultation Effect ($\Delta$DDI)}
The consultation effect measures post-consultation convergence:
\begin{equation}
\Delta\text{DDI} = \text{DDI}_A - \text{DDI}_C
\end{equation}
A positive $\Delta\text{DDI}$ indicates that parliamentary discourse moved closer to civic
input after the consultation closed; $\Delta\text{DDI} \approx 0$ indicates no detectable
effect. For cases where Phase~A is empty, $\Delta\text{DDI}$ is not computed.

\subsubsection{Statistical Validation}
\label{sec:stats}

To assess the reliability of DDI and $\Delta$DDI, we apply two complementary
statistical procedures.

\paragraph{Bootstrap confidence intervals.}
We resample Phase~B, Phase~A, and Phase~C independently with replacement ($n=2000$),
recomputing $d_{\text{salience}}$ and $d_{\text{stance}}$ from the resampled distributions;
$d_{\text{topic}}$ is held fixed (it depends on topic centroids, not individual documents).
Phase~B is resampled once per iteration and used jointly for both DDI(B,A) and DDI(B,C),
preserving their correlation in the $\Delta$DDI estimate.

\paragraph{Permutation test.}
We pool Phase~A and Phase~C speeches for each case, randomly split them into groups
of sizes $|A|$ and $|C|$ ($n=2000$ permutations), and compute the null distribution of $\Delta$DDI
under $H_0$: no temporal effect. The two-tailed $p$-value is the fraction of permuted
$|\Delta\text{DDI}|$ exceeding the observed $|\Delta\text{DDI}|$.

\paragraph{Cross-dossier null.}
We also compute $d_{\text{stance}}(B_X, C_Y)$ for all six cross-dossier pairings
($X \neq Y$) to verify that the metric assigns higher distance to unrelated corpus
pairs than to within-dossier pairs.

\subsubsection{Sensitivity Analysis}
To assess robustness, we re-compute DDI under three alternative weight configurations
(in addition to the base setting):
\begin{itemize}
  \item \textbf{Uniform}: $\mathbf{w} = (w_t, w_s, w_\sigma) = (0.33, 0.33, 0.34)$
  \item \textbf{Topic-heavy}: $\mathbf{w} = (0.60, 0.25, 0.15)$
  \item \textbf{Stance-heavy}: $\mathbf{w} = (0.25, 0.30, 0.45)$
\end{itemize}
If case rankings are invariant across configurations, we conclude the DDI is robust to
weight specification.

\section{Results}

\subsection{Topic Modeling}

BERTopic produced interpretable, language-diverse topic clusters in all phases.
Table~\ref{tab:topics} summarises the output per case and phase.

\begin{table}[h]
\centering
\small
\begin{tabular}{llrrrr}
\toprule
\textbf{Case} & & \textbf{Ph.~A} & \textbf{Ph.~B} & \textbf{Ph.~C} & \textbf{Unified} \\
\midrule
COVID Certificate  & Topics   & 17 & 98 & 7  & 72 \\
                   & Outlier\% & 11 & 33 & 7  & -- \\
Climate Target     & Topics   & 71 & 17 & 36 & 111 \\
                   & Outlier\% & 30 & 19 & 23 & -- \\
Pesticides         & Topics   & 16 & 75 & 15 & 67 \\
                   & Outlier\% & 20 & 41 & 17 & -- \\
\bottomrule
\end{tabular}
\caption{BERTopic output per phase. Outlier\% = fraction of documents assigned to topic $-1$.}
\label{tab:topics}
\end{table}

The high Phase~B outlier rate for Pesticides (41\%) reflects the heterogeneous multilingual
corpus (8,140 items, DE/PL/FR/EN); HDBSCAN does not force items into clusters, which is
preferable to inflating topic cohesion.

\subsection{Stance Detection}

Zero-shot NLI (mDeBERTa-v3, 24 languages, no translation) produces coherent distributions
across all phases. Table~\ref{tab:stance} reports the aggregate results for all three cases.

\begin{table}[h]
\centering
\small
\begin{tabular}{llrrr}
\toprule
\textbf{Case} & \textbf{Phase} & \textbf{Support (\%)} & \textbf{Neutral (\%)} & \textbf{Oppose (\%)} \\
\midrule
COVID Certificate & A (EP pre)   & 3.8  & 86.4 & 9.8  \\
                  & B (Citizens) & 54.2 & 44.9 & 0.9  \\
                  & C (EP post)  & 5.8  & 90.7 & 3.5  \\
\midrule
Pesticides SUR    & A (EP pre)   & 1.3  & 81.7 & 17.1 \\
                  & B (Citizens) & 2.0  & 48.1 & 49.9 \\
                  & C (EP post)  & 1.0  & 85.0 & 14.0 \\
\midrule
Climate Target    & A (EP pre)   & 1.9  & 96.8 & 1.4  \\
                  & B (Citizens) & 7.7  & 86.1 & 6.2  \\
                  & C (EP post)  & 1.9  & 96.0 & 2.2  \\
\bottomrule
\end{tabular}
\caption{Stance distributions per phase. Hypothesis: COVID = ``the EU digital COVID certificate should be extended''; Pesticides = ``the use of pesticides in EU agriculture should be reduced''; Climate Target = ``the EU should set ambitious climate targets for 2040.''}
\label{tab:stance}
\end{table}

\noindent
For \textbf{COVID}, citizens are predominantly \textit{supportive} of extension (54\%), while the EP
is largely neutral or procedural---a surprising reversal of the expected opposition signal,
likely due to the sample being pro-extension activists rather than opponents.
For \textbf{Pesticides}, citizens are evenly split between neutral (48\%) and oppose (50\%),
while the EP is overwhelmingly neutral (82--85\%) before and after---reflecting the technical
register of legislative debate vs.\ the polarised civic response.
For \textbf{Climate Target}, both EP phases are near-unanimously neutral (${\sim}96$--97\%),
while citizens show modest but symmetric activity (8\% support, 6\% oppose)---consistent
with the ideologically contested nature of the 2040 target, where no clear bloc dominates.

\subsection{Deliberative Distance Index}

Table~\ref{tab:ddi} reports DDI and $\Delta$DDI for all three cases.
Base weights: $w_{\text{topic}}=0.40$, $w_{\text{salience}}=0.35$, $w_{\text{stance}}=0.25$.
Components $d_t, d_s, d_\sigma$ are reported for the B\,vs.\,C comparison.

\begin{table*}[t]
\centering
\small
\begin{tabular}{lrrrrrrr}
\toprule
\textbf{Case} & $d_t$ & $d_s$ & $d_\sigma$ & \textbf{DDI(B,A)} & \textbf{DDI(B,C)} & $\Delta$\textbf{DDI} & \textbf{95\% CI} \\
\midrule
COVID Certificate  & 0.375 & 0.888 & 0.226 & 0.497 & 0.517 & $-$0.020 & [$-$0.049,~$+$0.006] \\
Climate Target     & 0.244 & 0.702 & 0.023 & 0.352 & 0.349 & $+$0.001 & [$-$0.013,~$+$0.015] \\
Pesticides SUR     & 0.294 & 0.792 & 0.116 & 0.397 & 0.424 & $-$0.027$^*$ & [$-$0.046,~$-$0.006] \\
\bottomrule
\end{tabular}
\caption{DDI components ($d_t$=topic, $d_s$=salience, $d_\sigma$=stance) and $\Delta$DDI
with bootstrap 95\% CI ($n=2000$). $\Delta\text{DDI} = \text{DDI}(B,A) - \text{DDI}(B,C)$;
positive values indicate post-consultation convergence.
$^*$CI excludes zero: statistically significant post-consultation divergence ($p=0.028$, permutation test).}
\label{tab:ddi}
\end{table*}

\textbf{Key findings}:
\begin{itemize}
  \item \textbf{Salience distance dominates all cases}: $d_s$ ranges from 0.70 to 0.89
        (log$_2$, $[0,1]$), making it the largest component of DDI. Citizens and the EP
        discuss many of the same themes (moderate $d_t \approx 0.24$--$0.38$) but weight
        them very differently. Stance divergence is smaller ($d_\sigma \leq 0.23$),
        indicating that the primary gap is \textit{what they emphasise}, not \textit{what they believe}.

  \item \textbf{COVID and Climate Target: null finding.}
        Bootstrap 95\% CIs for $\Delta$DDI include zero
        ([$-$0.049, $+$0.006] and [$-$0.013, $+$0.015] respectively),
        and permutation $p$-values are 0.72 and 0.64.
        No statistically detectable alignment effect is found for these two dossiers.

  \item \textbf{Pesticides SUR: statistically significant post-consultation divergence.}
        $\Delta\text{DDI} = -0.027$, 95\% CI [$-$0.046, $-$0.006] (excludes zero),
        permutation $p = 0.028$.
        The EP post-consultation corpus moved \emph{further} from citizen concerns
        rather than closer. This is consistent with the legislative outcome:
        the SUR regulation was ultimately withdrawn by the Commission in February 2024
        following opposition in the EP, suggesting a structural disconnect between
        citizen environmental advocacy and the parliamentary deliberation trajectory.

  \item \textbf{COVID} shows the highest absolute DDI ($\text{DDI}(B,C) = 0.517$),
        driven by near-maximal salience distance ($d_s = 0.888$).
        The wide bootstrap CI for $\Delta$DDI ([$-$0.049, $+$0.006]) reflects
        the small Phase~C corpus (86 speeches), limiting statistical power.

  \item \textbf{Climate Target} has the lowest DDI ($\approx 0.35$) and the
        tightest CI for $\Delta$DDI (width: 0.028), supported by the largest
        Phase~A and Phase~C corpora. The near-zero $\Delta$DDI is reliable
        and consistent with a mature dossier where EP and citizen corpora
        already share overlapping thematic structures.
\end{itemize}

\subsubsection{Sensitivity Analysis}
Table~\ref{tab:sensitivity} reports $\Delta$DDI across the four weight configurations
defined in Section~\ref{sec:ddi}.

\begin{table}[h]
\centering
\small
\begin{tabular}{lrrrr}
\toprule
\textbf{Case} & \textbf{Base} & \textbf{Uniform} & \textbf{Topic-heavy} & \textbf{Stance-heavy} \\
              & $(0.40,0.35,0.25)$ & $(0.33,0.33,0.34)$ & $(0.60,0.25,0.15)$ & $(0.25,0.30,0.45)$ \\
\midrule
COVID Certificate  & $-$0.020 & $-$0.013 & $-$0.027 & $-$0.004 \\
Climate Target     & $+$0.003 & $+$0.003 & $+$0.003 & $+$0.003 \\
Pesticides SUR     & $-$0.027 & $-$0.028 & $-$0.020 & $-$0.028 \\
\bottomrule
\end{tabular}
\caption{$\Delta$DDI sensitivity to weight configuration.
Positive = post-consultation convergence; negative = divergence.
Sign and rank are invariant across all configurations.}
\label{tab:sensitivity}
\end{table}

The case ranking is invariant: COVID and Pesticides show mild divergence ($\Delta$DDI $< 0$)
under all weight configurations; Climate Target shows near-zero positive $\Delta$DDI under all
configurations. The magnitude of COVID's $\Delta$DDI is most sensitive to the topic weight
(topic-heavy: $-0.025$; stance-heavy: $-0.006$), reflecting that $d_{\text{topic}}$ is the
component that differentiates pre- and post-consultation EP discourse for that case.

\subsubsection{Statistical Robustness}
\label{sec:robustness}

Table~\ref{tab:perm} reports the permutation test results and the cross-dossier null
baseline for $d_{\text{stance}}$.

\begin{table}[h]
\centering
\small
\begin{tabular}{lrrrl}
\toprule
\textbf{Case} & $\Delta$\textbf{DDI} & \textbf{Perm.\ null 95\% CI} & $p$\textbf{-value} & \textbf{Interpretation} \\
\midrule
COVID Certificate & $-$0.020 & [$-$0.062, $+$0.002] & 0.720 & Not significant \\
Climate Target    & $+$0.003 & [$-$0.012, $+$0.010] & 0.641 & Not significant \\
Pesticides SUR    & $-$0.027 & [$-$0.021, $+$0.026] & \textbf{0.028} & Significant divergence \\
\bottomrule
\end{tabular}
\caption{Permutation test results ($n=2000$). The permutation null pools Phase~A and Phase~C
speeches and randomly reassigns them to two groups of sizes $|A|$ and $|C|$.
$p$-value: two-tailed fraction of permuted $|\Delta\text{DDI}| \geq$ observed.}
\label{tab:perm}
\end{table}

\paragraph{Cross-dossier null.}
Mean within-dossier $d_{\text{stance}}$ is 0.121; mean cross-dossier
$d_{\text{stance}}(B_X, C_Y)$ for all six unrelated pairs is 0.191 (57\% higher).
This confirms that $d_{\text{stance}}$ assigns systematically lower distance to
corpus pairs from the same dossier, validating the metric's discriminative ability.
Individual cross-dossier values range from 0.004 (Climate$_B$--COVID$_C$,
both EP-neutral corpora) to 0.336 (COVID$_B$--Pesticides$_C$, opposing stance profiles).

\paragraph{Phase A robustness (Problem 3).}
Restricting Phase~A to speeches strictly before the consultation open date
(Phase~A\textsubscript{strict}) changes DDI(B,A) by at most 0.012 for COVID
and 0.004 for Climate; for Pesticides the overlap is 0.8\% and results are unchanged.
The qualitative conclusions are robust to the Phase~A temporal definition.

\subsection{Graph Metrics}

Table~\ref{tab:graph} reports topological metrics from the bipartite contribution--topic
graph (unified BERTopic model, B\,vs\,C subgraph).

\begin{table*}[t]
\centering
\small
\begin{tabular}{lrrrrr}
\toprule
\textbf{Case} & \textbf{CPBR} & \textbf{MPET} & \textbf{CPBC\textsubscript{max}} & $r_{BC}^{\text{all}}$ & $r_{BC}^{\text{bridge}}$ \\
\midrule
COVID Certificate & 0.088 & 0.324 & $6.2{\times}10^{-5}$ & $-$0.211 & $-$0.353$^\dagger$ \\
Climate Target    & 0.386 & 0.628 & $2.6{\times}10^{-4}$ & $-$0.098 & $-$0.179 \\
Pesticides SUR    & 0.333 & 0.272 & $7.9{\times}10^{-5}$ & $-$0.213 & $-$0.286 \\
\bottomrule
\end{tabular}
\caption{Graph metrics on the B--C bipartite contribution--topic subgraph.
CPBR = cross-phase bridging ratio; MPET = mean phase entropy (balance of shared topics between phases).
CPBC\textsubscript{max} = subset betweenness centrality of the top topic node on B$\to$C paths,
normalised by $|B|\times|C|$; small values ($\sim$$10^{-5}$--$10^{-4}$) are expected for sparse bipartite graphs
where most B--C paths pass through few shared hubs.
$r_{BC}^{\text{all}}$ = Pearson $r$ of per-phase degrees across all topic nodes (incl.\ zero-degree);
$r_{BC}^{\text{bridge}}$ = restricted to bridging topics ($n_{\text{bridge}}$=6, 39, 22 respectively).
$^\dagger$COVID $r_{BC}^{\text{bridge}}$ is unreliable ($n=6 < 10$).}
\label{tab:graph}
\end{table*}

\noindent
$r_{BC}^{\text{all}}$ is negative in all three cases, indicating that topics heavily used by
citizens (Phase~B) tend not to be the topics heavily used by the EP post-consultation (Phase~C).
$r_{BC}^{\text{bridge}}$, restricted to topics present in both phases, is more negative still,
showing that even among genuinely shared topics the relative intensities are anti-correlated.
This is a structural signature of agenda divergence fully consistent with the high $d_s$ values.
Note that COVID's $r_{BC}^{\text{bridge}}$ ($n=6$) is below the reliability threshold and should
not be interpreted.
Climate Target shows the highest CPBR (0.386) and MPET (0.628) and the least negative
$r_{BC}^{\text{all}}$ ($-$0.098), consistent with its lowest absolute DDI ($\approx 0.35$): more
topic nodes bridge both phases, and shared topics are more evenly used. COVID has the lowest
CPBR (0.088), corroborating its highest DDI (0.517) and near-maximal salience divergence.
Pesticides occupies the middle position: moderate CPBR (0.333) but low MPET (0.272),
indicating shared topics are dominated by one phase---consistent with its intermediate DDI (0.42).

\section{Discussion}

\subsection{Implications for EU Better Regulation}

\subsubsection{Differentiated Findings Across Cases}
Statistical analysis reveals a more nuanced picture than a uniform null result.
For \textbf{COVID} and \textbf{Climate Target}, $\Delta$DDI CIs include zero (permutation $p > 0.6$):
no detectable alignment effect, consistent with \citet{Neumann2025}'s finding that
HYS is dominated by organised stakeholders and \citet{Bunea2024}'s observation that
consultation policy spaces are structurally low-dimensional.
For \textbf{Pesticides}, $\Delta$DDI $= -0.027$ is statistically significant ($p = 0.028$),
indicating post-consultation \emph{divergence}: the EP moved further from citizen concerns
after the consultation closed. This finding is consistent with the legislative trajectory:
the Sustainable Use Regulation was ultimately withdrawn by the Commission in February~2024
following EP opposition, signalling a structural disconnect between the citizen
advocacy corpus (50\% opposition stance against the EP's predominantly neutral register)
and the parliamentary deliberation path.
This illustrates the DDI's utility as a monitoring tool: the Pesticides case would
have been flagged at the time of consultation as a high-divergence, high-salience-gap
dossier warranting further legislative attention.

\subsubsection{Consultation Design and Timing}
COVID's null finding is consistent with a post-hoc consultation that opened after
the main EP debate was already underway (Phase~A covers July 2021 -- May 2022,
predating the extension proposal of January 2022 by several months).
Climate Target's null finding may reflect a mature policy space where
EP and citizen discourse had already converged on shared vocabulary (lowest $d_t = 0.24$)
before the consultation opened.

\subsubsection{Salience as Primary Signal}
Across all three cases, topic and stance distances are moderate and comparable.
Salience distance---how much weight each phase places on shared topics---is the dominant
discriminator. This suggests that future DDI applications should weight $d_s$ more heavily
when comparing forums with similar vocabularies but different thematic emphasis.

\subsubsection{Topic Gap as Legitimacy Indicator}
Topics present in Phase~B (citizen corpus) but absent in Phase~C (post-consultation EP) represent
citizen concerns unaddressed in the legislative debate. Identifying these gaps computationally
is a key practical use of the DKG infrastructure.

\subsection{Theoretical Contributions}

\subsubsection{Operationalizing Deliberative Quality}
DDI translates normative ideals (inclusiveness, responsiveness) into measurable metrics. Unlike manual deliberative quality coding \citep{Steenbergen2003}, DDI scales to thousands of cases.

\subsubsection{Multi-Forum Deliberation}
Extending deliberative democracy theory \citep{Dryzek2000} to fragmented digital spaces. The DKG enables systemic analysis of the \textit{deliberative system} \citep{Parkinson2012}, not isolated forums.

\subsection{Comparison to Di Porto et al. (2024)}

Where \citet{DiPorto2024} apply NLP clustering to 830 HYS responses to extract topic lists (single-forum, descriptive), our approach links HYS to the EP via BERTopic alignment and stance detection to compute a cross-forum distance metric. The key advance is moving from topic mining to \emph{cross-forum alignment measurement} over 13k+ contributions and 4k+ speeches, enabling auditing rather than description.

\section{Limitations and Future Work}

\subsection{Current Limitations}
\begin{itemize}
  \item \textbf{EP coverage}: Plenary debates only; excludes committee hearings and trilogue negotiations, where the legislative detail is actually worked out.

  \item \textbf{Stance detection register bias}: The zero-shot NLI model (mDeBERTa-v3) classifies EP speeches predominantly as ``neutral'' (86--97\%), partly because parliamentary language is procedural and hedged rather than declarative. This conflates genuine policy neutrality with technical register, inflating $d_{\text{stance}}$ between citizen and EP corpora. Mitigation: fine-tuning on parliamentary text, or restricting to the most argumentative speech segments, is left for future work.

  \item \textbf{Topic vocabulary not grounded in EuroVoc}: BERTopic produces data-driven, case-specific topic clusters whose labels are not aligned with controlled vocabulary. This limits cross-dossier interpretability. Mapping BERTopic centroids to EuroVoc descriptors (via nearest-neighbour in the shared MiniLM space) would provide EU-standard thematic labels and enable comparison across dossiers; this is a planned extension.

  \item \textbf{Corpus size asymmetry}: COVID Phase~C (86 speeches, 7 topics) limits the statistical power of $\text{DDI}(B, C)$ for that case. The bootstrap CI is correspondingly wide ([$-$0.049, $+$0.006]).

  \item \textbf{Phase A temporal bounds}: The current implementation includes speeches made during the consultation window in Phase~A (up to 37\% for Climate Target). A stricter cutoff at $t_{\text{open}}$ reduces DDI(B,A) by at most 0.012; conclusions are robust (Section~\ref{sec:robustness}).

  \item \textbf{DDI weights}: Weights $(0.40, 0.35, 0.25)$ are set a priori; sensitivity analysis across four configurations shows the Pesticides result is robust ($p \leq 0.05$ under all configurations) \citep{OECD2008}.

  \item \textbf{Causal claims}: $\Delta$DDI measures association; consultations may not causally shift EP agenda. Instrumental-variable designs are needed for causal attribution.

  \item \textbf{Generalizability}: Three cases cannot support EU-wide claims; expansion to 20+ dossiers with ground-truth alignment labels is needed for DDI calibration.
\end{itemize}

\subsection{Future Extensions}

\begin{itemize}
  \item \textbf{EuroVoc grounding}: Map BERTopic centroids to nearest EuroVoc descriptor in the shared MiniLM space, enabling cross-dossier comparison by EU thematic domain.
  \item \textbf{Long-document stance}: Sliding-window chunk-and-aggregate to reduce 512-token truncation bias in long HYS submissions.
  \item \textbf{Cross-national extension}: Apply DDI to national parliaments via the ParlaMint corpus \citep{Erjavec2023} and compare deliberative distance across EU member states.
  \item \textbf{Real-time auditing}: Deploy DDI as a live monitoring dashboard to flag divergence during active consultations.
  \item \textbf{Causal inference}: Use instrumental-variable designs to test whether consultations causally shift EP agenda, rather than measuring association only.
\end{itemize}

\section{FAIR Principles and Sustainability}

\subsection{FAIR Assessment}
The DKG is designed to satisfy the FAIR data principles \citep{Wilkinson2016}:

\begin{itemize}
  \item \textbf{Findable}: The ontology is served at a persistent URI (\url{https://w3id.org/deliberation/ontology}) registered via w3id.org. The code repository is indexed on GitHub. Dataset metadata follows DCAT-2 vocabulary.
  \item \textbf{Accessible}: All data are available as Linked Open Data via the SPARQL endpoint (\url{https://dkg.svagnoni.linkeddata.es/sparql}), as downloadable Turtle (TTL) dumps, and via a human-readable web interface. No authentication is required for read access.
  \item \textbf{Interoperable}: The ontology reuses established vocabularies (\texttt{rdf}, \texttt{rdfs}, \texttt{owl}, \texttt{xsd}, \texttt{dct}, \texttt{prov-o}) and aligns with EuroVoc and SKOS for thematic annotation. RDF triples link HYS contributions to EP debates via shared EUR-Lex procedure identifiers (CELEX/Ares).
  \item \textbf{Reusable}: All data and derived outputs are released under CC~BY~4.0. The ontology includes \texttt{rdfs:label} and \texttt{rdfs:comment} annotations in English. A SPARQL endpoint, a web interface, and a citizen-facing semantic search application (\url{https://citizen.svagnoni.linkeddata.es}) support diverse reuse scenarios.
\end{itemize}

\subsection{Sustainability Plan}
The DKG is maintained as an ongoing research infrastructure with the following commitments:
\begin{itemize}
  \item \textbf{Ontology versioning}: The \texttt{del:} ontology is version-controlled in the public repository; breaking changes will be signalled via new version IRIs following the w3id.org pattern.
  \item \textbf{Data updates}: The HYS scraper runs periodically to update the feedback corpus; EP debate ingestion targets quarterly updates aligned with plenary session schedules.
  \item \textbf{Endpoint uptime}: The Apache Jena Fuseki triple store and the citizen interface are hosted on a dedicated server; service continuity is maintained by the research group.
  \item \textbf{Community reuse}: The pilot study dataset and analysis scripts are designed for extension; we provide a documented API and example SPARQL queries to lower the barrier for new contributors.
\end{itemize}

\section{Reproducibility and Open Science}

\subsection{Data Availability}
\begin{itemize}
  \item \textbf{Pilot dataset}: available in the DKG repository under \texttt{data/pilot\_study\_dataset\_v2/} (CC BY 4.0)
  \item \textbf{DKG SPARQL endpoint}: \url{https://dkg.svagnoni.linkeddata.es/sparql}
  \item \textbf{DKG website and explorer}: \url{https://dkg.svagnoni.linkeddata.es}
  \item \textbf{Ontology}: \url{https://w3id.org/deliberation/ontology} (persistent URI via w3id.org)
  \item \textbf{Source databases}: HYS feedback scraped from EC Better Regulation API; EP verbatim from EP Open Data Portal
\end{itemize}

\subsection{Code Availability}
\begin{itemize}
  \item \textbf{Repository}: \url{https://github.com/Stocastico96/Deliberation-Knowledge-Graph}
  \item \textbf{Analysis scripts}: \texttt{scripts/analysis/00\_export\_cases.py} through \texttt{05c\_null\_baseline.py} (including bootstrap CI and permutation null)
  \item \textbf{Dependencies}: BERTopic 0.16+, sentence-transformers 2.2+, transformers 4.30+, networkx 3.0+
  \item \textbf{Compute requirements}: 16 GB RAM, 4-core CPU (no GPU required for pilot)
  \item \textbf{License}: MIT (code); CC BY 4.0 (data and derived outputs)
\end{itemize}

\subsection{Reproducibility Instructions}
\begin{lstlisting}[language=bash, basicstyle=\small\ttfamily]
git clone https://github.com/Stocastico96/Deliberation-Knowledge-Graph
cd Deliberation-Knowledge-Graph
pip install -r requirements.txt

# Step 1a: per-phase BERTopic models (for d_topic)
python3 scripts/analysis/01_bertopic_modeling.py
# Step 1b: unified BERTopic model (for d_salience)
python3 scripts/analysis/01b_unified_topic_model.py

# Step 2: cross-forum topic alignment
python3 scripts/analysis/02_topic_alignment.py

# Step 3: zero-shot stance detection
python3 scripts/analysis/03_stance_detection.py

# Step 4a: DDI computation
python3 scripts/analysis/04_calculate_ddi.py
# Step 4b: topological graph metrics
python3 scripts/analysis/04b_graph_metrics.py

# Step 5b: bootstrap confidence intervals on DDI and delta-DDI
python3 scripts/analysis/05b_bootstrap_ci.py
# Step 5c: permutation null and cross-dossier null baseline
python3 scripts/analysis/05c_null_baseline.py
\end{lstlisting}

\section{Conclusions}

This paper applied the Deliberation Knowledge Graph to quantify deliberative distance between EU citizen consultations and parliamentary debates.
The Deliberative Distance Index reveals that \textit{salience divergence} is the primary driver of cross-forum distance: citizens and the EP share overlapping topic vocabularies but weight themes very differently ($d_s \approx 0.70$--$0.89$).
Bootstrap and permutation tests yield a differentiated picture.
For COVID and Climate Target, $\Delta$DDI is not statistically different from zero
($p > 0.6$): no detectable alignment effect within the observation window.
For Pesticides, $\Delta\text{DDI} = -0.027$ is statistically significant ($p = 0.028$,
bootstrap CI $[-0.046, -0.006]$), indicating post-consultation \emph{divergence}---consistent with the SUR regulation's withdrawal in February 2024.
Topological graph metrics corroborate the cross-sectional ranking: Climate Target has the lowest DDI and the highest CPBR and MPET; COVID has the highest DDI and the lowest CPBR and most negative $r_{BC}$.
A cross-dossier null baseline confirms that $d_{\text{stance}}$ is 57\% higher for unrelated corpus pairs than for same-dossier pairs, validating the metric's discriminative ability.

The DKG provides a scalable, reproducible framework for auditing deliberative alignment in EU lawmaking and beyond. By integrating heterogeneous data sources under a unified ontology and exposing them via a public SPARQL endpoint, it enables cross-platform, cross-national, and longitudinal research on democratic deliberation. All data, code, and query endpoints are openly available to support replication and extension.

\textbf{Key takeaway}: The pilot study reveals a structural salience gap between citizen and parliamentary discourse. For two of three cases (COVID, Climate Target), $\Delta$DDI is not statistically distinguishable from zero, a meaningful null result suggesting that deliberative alignment cannot be assumed to arise automatically from the Better Regulation process. For Pesticides, the statistically significant post-consultation divergence ($p=0.028$) demonstrates that the DDI can detect adverse alignment trajectories: the metric would have flagged this dossier as high-divergence at the time of consultation, consistent with the regulation's subsequent withdrawal. This dual finding motivates monitoring tools such as the DDI to surface alignment gaps in real time.

\section*{Resource Availability Statement}
\noindent
The Deliberation Knowledge Graph is publicly available at:
\begin{itemize}
  \item \textbf{Ontology (persistent URI)}: \url{https://w3id.org/deliberation/ontology} (CC BY 4.0)
  \item \textbf{SPARQL endpoint}: \url{https://dkg.svagnoni.linkeddata.es/sparql} (open, no authentication)
  \item \textbf{Web interface}: \url{https://dkg.svagnoni.linkeddata.es}
  \item \textbf{Code repository}: \url{https://github.com/Stocastico96/Deliberation-Knowledge-Graph} (MIT licence)
  \item \textbf{Pilot dataset}: available in the repository under \texttt{data/pilot\_study\_dataset\_v2/} (CC BY 4.0)
  \item \textbf{Canonical citation}: \citet{Vagnoni2026}
\end{itemize}
All resources are accessible without restriction and will remain available for the duration of the review process and beyond.

\section*{Acknowledgments}
We thank the European Commission for open access to Have Your Say data, and the European Parliament for publicly available verbatim transcripts. This research received no specific grant funding.

\begin{thebibliography}{99}

\bibitem[Vagnoni \& Rodr\'iguez-Doncel(2026)]{Vagnoni2026}
Vagnoni, S., \& Rodr\'iguez-Doncel, V. (2026).
\textit{A Deliberation Knowledge Graph: Bridging Institutional and Civic Democratic Discourse}.
In \textit{Electronic Government and the Information Systems Perspective}, pp. 122--136.
Springer Nature Switzerland.
doi: 10.1007/978-3-032-02225-7_9

\bibitem[Di Porto et al.(2024)]{DiPorto2024}
Di Porto, F., Maggiolino, M., \& Parcu, P.L. (2024).
Mining EU consultations through AI.
\textit{Artificial Intelligence and Law}, 1--38.

\bibitem[EC(2015)]{EC2015}
European Commission (2015).
Better Regulation Guidelines.
SWD(2015) 111 final.

\bibitem[Abercrombie \& Batista-Navarro(2020)]{Abercrombie2020}
Abercrombie, G., \& Batista-Navarro, R. (2020).
Sentiment and position-taking analysis of parliamentary debates: A systematic literature review.
\textit{Journal of Computational Social Science}, 3, 245--270.

\bibitem[Hix et al.(2006)]{Hix2006}
Hix, S., Noury, A., \& Roland, G. (2006).
Dimensions of politics in the European Parliament.
\textit{American Journal of Political Science}, 50(2), 494--520.

\bibitem[Steenbergen et al.(2003)]{Steenbergen2003}
Steenbergen, M.R., B\"achtiger, A., Sp\"orndli, M., \& Steiner, J. (2003).
Measuring political deliberation: A Discourse Quality Index.
\textit{Comparative European Politics}, 1(1), 21--48.

\bibitem[Lawrence \& Reed(2020)]{Lawrence2020}
Lawrence, J., \& Reed, C. (2020).
Argument mining: A survey.
\textit{Computational Linguistics}, 45(4), 765--818.

\bibitem[Palmirani et al.(2018)]{Palmirani2018}
Palmirani, M., Martoni, M., Rossi, A., Bartolini, C., \& Robaldo, L. (2018).
Legal ontology for modelling GDPR concepts and norms.
\textit{Legal Knowledge and Information Systems}, 91--100.

\bibitem[Mader \& Haslhofer(2012)]{Mader2012}
Mader, C., \& Haslhofer, B. (2012).
Perception and relevance of linked open data.
\textit{International Journal on Semantic Web and Information Systems}, 8(3), 1--24.

\bibitem[Dryzek(2000)]{Dryzek2000}
Dryzek, J.S. (2000).
\textit{Deliberative Democracy and Beyond: Liberals, Critics, Contestations}.
Oxford University Press.

\bibitem[Parkinson \& Mansbridge(2012)]{Parkinson2012}
Parkinson, J., \& Mansbridge, J. (Eds.) (2012).
\textit{Deliberative Systems: Deliberative Democracy at the Large Scale}.
Cambridge University Press.

\bibitem[Erjavec et al.(2023)]{Erjavec2023}
Erjavec, T., et al. (2023).
The ParlaMint corpora of parliamentary proceedings.
\textit{Language Resources and Evaluation}, 57, 415--448.

\bibitem[Reimers \& Gurevych(2019)]{Reimers2019}
Reimers, N., \& Gurevych, I. (2019).
Sentence-BERT: Sentence embeddings using Siamese BERT-networks.
\textit{Proceedings of EMNLP 2019}, 3982--3992.

\bibitem[Adam \& Kogler(2024)]{Adam2024}
Adam, R., \& Kogler, M.~L. (2024).
Bidirectional Topic Matching: Quantifying thematic overlap between corpora through topic modelling.
arXiv preprint arXiv:2412.18376.

\bibitem[Lu et al.(2020)]{Lu2020}
Lu, J., Henchion, M., \& Mac~Namee, B. (2020).
Diverging divergences: Examining variants of Jensen Shannon divergence for corpus comparison tasks.
In \textit{Proceedings of LREC 2020}, 6740--6744.

\bibitem[Bunea et al.(2024)]{Bunea2024}
Bunea, A., W\"uest, R., \& Lipcean, S. (2024).
Mapping the policy space of public consultations: Evidence from the European Union.
\textit{Journal of European Public Policy}, 32(3), 755--783.

\bibitem[Neumann et al.(2025)]{Neumann2025}
Neumann, R., Lang, S., \& Meng, P. (2025).
`Have Your Say' in practice: Assessing citizens' use of the EU's public consultation platform.
\textit{European Politics and Society}, 26(5), 1122--1142.

\bibitem[Lin(1991)]{Lin1991}
Lin, J. (1991).
Divergence measures based on the Shannon entropy.
\textit{IEEE Transactions on Information Theory}, 37(1), 145--151.

\bibitem[OECD(2008)]{OECD2008}
Nardo, M., Saisana, M., Saltelli, A., Tarantola, S., Hoffmann, A., \& Giovannini, E. (2008).
\textit{Handbook on Constructing Composite Indicators: Methodology and User Guide}.
OECD Publishing.

\bibitem[Wilkinson et al.(2016)]{Wilkinson2016}
Wilkinson, M.D., et al. (2016).
The FAIR Guiding Principles for scientific data management and stewardship.
\textit{Scientific Data}, 3, 160018.

\end{thebibliography}

\section*{Declaration of Use of Generative AI}

In preparing this manuscript, the authors used Claude Sonnet (Anthropic) to assist with:
drafting and editing portions of the text; identifying inconsistencies between the code and the paper;
and verifying that cited claims matched source materials.
All scientific results, experiment designs, code implementations, and final editorial decisions were made by the human authors, who take full responsibility for the work.
No generative AI was used to produce figures, compute results, or generate citations autonomously.

\end{document}
